\chapter{1999:Ahmed Zewail}

\label{chap:chemistry1999}

The Nobel Prize in Chemistry for the year 1999 was awarded to Ahmed Zewail.

\textbf{Motivation:}

for his studies of the transition states of chemical reactions using femtosecond spectroscopy

\section{Ahmed Zewail}

\subsection*{Early Life and Education}

Ahmed Zewail was born in 1946 in Damanhur, Egypt.

\subsection*{Career and Major Research}

Zewail studied at Alexandria University (B.S., M.S.) and earned his Ph.D. from the University of Pennsylvania in 1974. He joined the California Institute of Technology in 1976, where he became the Linus Pauling Professor of Chemistry and Physics.

\subsection*{Nobel Prize-winning Work}

The work for which Ahmed Zewail was awarded the Nobel Prize in Chemistry in 1999 was described by the Nobel Committee as: "for his studies of the transition states of chemical reactions using femtosecond spectroscopy".

He developed femtosecond spectroscopy, using laser pulses of a few femtoseconds to observe the breaking and forming of chemical bonds in real time, thereby directly studying the transition states of reactions.

\subsection*{Significance and Impact}

The new field of femtochemistry allowed the first direct observation of the short-lived transition states of chemical reactions.

\subsection*{Later Career and Honors}

Zewail remained at the California Institute of Technology. He received the 1999 Nobel Prize in Chemistry, and he died in 2016.

\subsection*{Legacy}

Femtochemistry opened a window onto the fastest steps of chemical change and influenced fields from molecular biology to materials science.

\subsection*{Modern Developments}

Zewail's femtosecond spectroscopy, which he called femtochemistry, enabled scientists to observe chemical reactions in real time as bonds break and form. His techniques are now used to study photosynthesis, vision, and DNA photochemistry, and femtosecond lasers underpin modern microscopy, materials science, and manufacturing.

\subsection*{Selected Publications}

"Femtochemistry: Atomic-Scale Dynamics of the Chemical Bond" (Journal of Physical Chemistry, 1993)

\clearpage

\section*{Chapter Summary}

The 1999 Nobel Prize in Chemistry recognized the study of the transition states of chemical reactions. Zewail's femtosecond spectroscopy made it possible to observe chemical bond making and breaking in real time.
